\chapter{Conservation Laws}

\textbooknumber{\S 7}
\section{Momentum}

% ============ Problem 1 ============ %

\begin{problem}[title={Refraction of a particle at a potential step}]
{
A particle of mass $m$ moving with velocity $\mathbf{v}_1$ leaves a half-space in which its potential energy is a constant $U_1$ and enters another in which its potential energy is a different constant $U_2$. Determine the change in the direction of motion of the particle.
}
{
The potential energy only depend on the co-ordinate perpendicular to the plane separating the half-space. Thus, the component of the momentum in that plane is conserved. Denoting by $\theta_1$ and $\theta_2$ the angles between the normal to the plane and the velocities $v_1$ and $v_2$ of the particle before and after passing the plane, we have
\begin{align*}
    & P_1 \sin{\theta_1} = P_2 \sin{\theta_2} \\
    \Rightarrow & v_1 \sin{\theta_1} = v_2 \sin{\theta_2} \\
    \Rightarrow & \frac{\sin{\theta_1}}{\sin{\theta_2}} = \frac{v_2}{v_1} . \numberthis \label{S7P1_P}
\end{align*}
The potential energy of the system is also independent of time, therefore the energy of the particle is conserved. Posing $E_1 = T_1 + U_1$ and $E_2 = T_2 + U_2$ as the energy of the particle before and after passing the plane, the law of conservation of energy requires
\begin{align*}
    &E_1 = E_2 \\
    \Rightarrow & T_1 + U_1 = T_2 + U_2 \\
    \Rightarrow &\frac{1}{2} m v_1^2 + U_1 = \frac{1}{2} m v_2^2 + U_2 \\
    \Rightarrow & v_2^2 = v_1^2 + \frac{2}{m}(U_1-U_2). \numberthis \label{S7P1_E}
\end{align*}
By substituting \eqref{S7P1_E} in the square of \eqref{S7P1_P}, we get
\begin{equation*}
    \left( \frac{\sin{\theta_1}}{\sin{\theta_2}} \right)^2 = \frac{v_1^2 + \frac{2}{m}(U_1-U_2)}{v_1^2}.
\end{equation*}
After taking the square root, the result is  
}
{
\begin{equation*}
    \frac{\sin{\theta_1}}{\sin{\theta_2}} = \sqrt{1 + \frac{2}{mv_1^2}(U_1-U_2)}
\end{equation*}
}
[
We check selected cases by integrating Newton's equations for a particle with $m=1$~kg and $v_1=2$~m/s in the smooth potential $U(x)=\Delta U[1+\tanh(x/w)]/2$, with the consistent force $F_x=-\Delta U\,\operatorname{sech}^2(x/w)/(2w)$ and $F_y=0$. Starting at $x=-1$~m, we measure the outgoing velocity after reaching $x=1$~m or returning to $x=-1$~m. For $\Delta U=1$~J, incidence angles $0^\circ,15^\circ,30^\circ,40^\circ,44^\circ,44.5^\circ$ transmit and $45.5^\circ,46^\circ,50^\circ,60^\circ$ reflect, as predicted by $mv_1^2\cos^2\theta_1/2<\Delta U$ for reflection. The exact threshold $45^\circ$ is excluded. We also test a downward step $\Delta U=-0.6$~J at $0^\circ,30^\circ,60^\circ$. Reducing $w$ from $0.2$ to $0.1$ to $0.05$~m removes finite-endpoint tail effects; independently halving the RK4 step from $0.008$ to $0.001$~s reduces integration errors, checked against energy conservation with the actual endpoint potentials. At the finest settings, the angle error relative to the half-space formula is below $10^{-9}$~rad, tangential momentum is constant to $10^{-12}$~kg\,m/s, and the energy-balance error divided by the initial energy is below $10^{-9}$. These checks support the result in the selected cases; they are not a general proof.
\begin{figure}[H]
    \centering
    \includegraphics{Figures/S7P1-trajectories.pdf}
    \caption{Engine trajectories at the upward step: transmission at $30^\circ$ and reflection at $60^\circ$, with $w=0.05$~m. The dashed line is the interface $x=0$.}
    \label{fig:S7P1-trajectories}
\end{figure}
\begin{figure}[H]
    \centering
    \includegraphics{Figures/S7P1-angles.pdf}
    \caption{Predicted transmission angles (curves) and measured outgoing angles (points). The vertical line marks the upward step's reflection threshold, which is not sampled.}
    \label{fig:S7P1-angles}
\end{figure}
]
\end{problem}

\textbooknumber{\S 8}
\section{Centre of mass}

% ============ Problem 1 ============ %

\begin{problem}[title={Action transformed between inertial frames}]
{
Find the law of transformation of the action $S$ from one inertial frame to another.
}
{
The Lagrangian $L$ and $L'$ of a mechanical system in two inertial frames of reference $K$ and $K'$ are respectively
\begin{equation*}
    L = T - U = \frac{1}{2} \sum_a m_a v_a^2 - U
\end{equation*}
and
\begin{equation*}
    L' = T' - U = \frac{1}{2} \sum_a m_a v_a^{'2} - U.
\end{equation*}
If the frame $K'$ moves with velocity $\mathbf{V}$ relative to the frame $K$, the velocities of the particles of the mechanical system relative to the two frames are related by $\mathbf{v}_a = \mathbf{v}_a' + \mathbf{V}$. We can now express the relation of the Lagrangian of the system in the two frames by
\begin{align*}
    L &= \frac{1}{2} \sum_a m_a (\mathbf{v}_a' + \mathbf{V})^2 - U \\ 
    &= \frac{1}{2} V^2 \sum_a m_a + \mathbf{V} \cdot \sum_a m_a \mathbf{v}_a' + \frac{1}{2} \sum_a m_a v_a^{'2} - U \\
    &= L' + \mathbf{V} \cdot \mathbf{P}' + \frac{1}{2} \mu V^2 . \numberthis \label{S8P1_L}
\end{align*}
Integrating \eqref{S8P1_L} with respect to time, we obtain
\begin{align*}
    S &= \int L dt \\
    &= S' + \int \left( \mathbf{V} \cdot \mathbf{P}' \right)dt + \frac{1}{2} \mu V^2 t \\
    &= S' + \mathbf{V} \cdot \int \sum_a m_a \mathbf{v}_a' dt + \frac{1}{2} \mu V^2 t \\
    &= S' + \mathbf{V} \cdot \sum_a m_a \mathbf{r}_a' + \frac{1}{2} \mu V^2 t .
\end{align*}
The law of transformation of the action $S$ is then
}
{
\begin{equation*}
    S = S' + \mu \mathbf{V} \cdot \mathbf{R}' + \frac{1}{2} \mu V^2 t
\end{equation*}
}
\end{problem}

\textbooknumber{\S 9}
\section{Angular momentum}

% ============ Problem 1 ============ %

\begin{problem}[title={Angular momentum in cylindrical co-ordinates}]
{
Obtain the expressions for the Cartesian components and the magnitude of the angular momentum of a particle in cylindrical co-ordinates $r,\phi,z$.
}
{
The Cartesian components of the angular momentum are simply
\begin{equation*}
    \mathbf{M} = \mathbf{r} \times \mathbf{p} = \begin{bmatrix} x \\ y \\ z \end{bmatrix} \times m \begin{bmatrix} \Dot{x} \\ \Dot{y} \\ \Dot{z} \end{bmatrix} = m \begin{bmatrix} y \Dot{z} - z \Dot{y} \\ z \Dot{x} - x \Dot{z} \\ x \Dot{y} - y \Dot{x} \end{bmatrix} . \numberthis \label{S9P1_M}
\end{equation*}
In cylindrical co-ordinates, the Cartesian components are expressed by
\begin{equation*}
    \begin{bmatrix} x \\ y \\ z \end{bmatrix} = \begin{bmatrix} r \cos{\phi} \\ r \sin{\phi} \\ z \end{bmatrix}
\end{equation*}
and
\begin{equation*}
    \begin{bmatrix} \Dot{x} \\ \Dot{y} \\ \Dot{z} \end{bmatrix} = \begin{bmatrix} \Dot{r}\cos{\phi} - r\Dot{\phi}\sin{\phi} \\ \Dot{r}\sin{\phi} + r\Dot{\phi}\cos{\phi} \\ \Dot{z} \end{bmatrix}.
\end{equation*}
By substituting those in \eqref{S9P1_M}, we get
\begin{align*}
    \mathbf{M} &= m \begin{bmatrix} y \Dot{z} - z \Dot{y} \\ z \Dot{x} - x \Dot{z} \\ x \Dot{y} - y \Dot{x} \end{bmatrix} \\
    &= m \begin{bmatrix} (r\Dot{z} - z\Dot{r})\sin{\phi} - rz\Dot{\phi}\cos{\phi} \\ -(r\Dot{z} - z\Dot{r})\cos{\phi} - rz\Dot{\phi}\sin{\phi} \\ r^2\Dot{\phi} \end{bmatrix} . \numberthis \label{S9P1_M2}
\end{align*}
By taking the magnitude of \eqref{S9P1_M2}, we finally get
}
{
\begin{align*}
    M_x &= m(r\Dot{z} - z\Dot{r})\sin{\phi} - mrz\Dot{\phi}\cos{\phi} \\
    M_y &= -m(r\Dot{z} - z\Dot{r})\cos{\phi} - mrz\Dot{\phi}\sin{\phi} \\
    M_z &= mr^2\Dot{\phi} \\
    M^2 &= m^2r^2\Dot{\phi}^2(r^2 + z^2) + m^2(r\Dot{z} - z\Dot{r})^2
\end{align*}
}
\end{problem}

% ============ Problem 2 ============ %

\begin{problem}[title={Angular momentum in spherical co-ordinates}]
{
Obtain the expressions for the Cartesian components and the magnitude of the angular momentum of a particle in spherical co-ordinates $r,\theta,\phi$.
}
{
Let $\theta$ be the polar angle (measured from the $z$ axis) and $\phi$ the azimuth (measured in the $xy$ plane from the $x$ axis). In spherical co-ordinates, the Cartesian components are expressed by
\begin{equation*}
    \begin{bmatrix} x \\ y \\ z \end{bmatrix} = \begin{bmatrix} r \sin{\theta} \cos{\phi} \\ r \sin{\theta} \sin{\phi} \\ r \cos{\theta} \end{bmatrix}
\end{equation*}
and
\begin{equation*}
    \begin{bmatrix} \Dot{x} \\ \Dot{y} \\ \Dot{z} \end{bmatrix} = \begin{bmatrix} \Dot{r}\sin{\theta}\cos{\phi} + r\Dot{\theta}\cos{\theta}\cos{\phi} - r\Dot{\phi}\sin{\theta}\sin{\phi} \\ \Dot{r}\sin{\theta}\sin{\phi} + r\Dot{\theta}\cos{\theta}\sin{\phi} + r\Dot{\phi}\sin{\theta}\cos{\phi} \\ \Dot{r}\cos{\theta} - r\Dot{\theta}\sin{\theta} \end{bmatrix}.
\end{equation*}
By substituting those in \eqref{S9P1_M} of \probref{\S 9}{1}, we get
\begin{align*}
    \mathbf{M} &= m \begin{bmatrix} y \Dot{z} - z \Dot{y} \\ z \Dot{x} - x \Dot{z} \\ x \Dot{y} - y \Dot{x} \end{bmatrix} \\
    &= m \begin{bmatrix} -r^2(\Dot{\theta}\sin{\phi} + \Dot{\phi}\sin{\theta}\cos{\theta}\cos{\phi}) \\ r^2(\Dot{\theta}\cos{\phi} - \Dot{\phi}\sin{\theta}\cos{\theta}\sin{\phi}) \\ r^2\Dot{\phi}\sin^2{\theta} \end{bmatrix} . \numberthis \label{S9P2_M}
\end{align*}
By taking the magnitude of \eqref{S9P2_M}, we finally get
}
{
\begin{align*}
    M_x &= -mr^2(\Dot{\theta}\sin{\phi} + \Dot{\phi}\sin{\theta}\cos{\theta}\cos{\phi}) \\
    M_y &= mr^2(\Dot{\theta}\cos{\phi} - \Dot{\phi}\sin{\theta}\cos{\theta}\sin{\phi}) \\
    M_z &= mr^2\Dot{\phi}\sin^2{\theta} \\
    M^2 &= m^2r^4(\Dot{\theta}^2 + \Dot{\phi}^2\sin^2{\theta})
\end{align*}
}
\end{problem}

% ============ Problem 3 ============ %

\begin{problem}[title={Conserved momenta in symmetric fields}]
{
Which components of momentum $\mathbf{P}$ and angular momentum $\mathbf{M}$ are conserved in motion in the following fields?
}{}{}
\end{problem}
\begin{subproblem}
{
The field of an infinite homogeneous plane.
}
{
Let us first see how a symmetry of the field leads to the conservation of a component. For a particle of mass $m$ in a stationary external field $U$, the Lagrangian in Cartesian co-ordinates is
\begin{equation*}
    L = \frac{1}{2} m (\Dot{x}^2 + \Dot{y}^2 + \Dot{z}^2) - U(x,y,z).
\end{equation*}
Since $P_x = \pd{L}{\Dot{x}} = m\Dot{x}$, Lagrange's equation for $x$ gives
\begin{equation*}
    \Dot{P}_x = \frac{d}{dt} \pd{L}{\Dot{x}} = \pd{L}{x} = -\pd{U}{x} .
\end{equation*}
If the field is unchanged by any translation along $x$, \ie $U(x+a,y,z) = U(x,y,z)$ for every $a$, then $U$ does not depend on $x$ and $P_x$ is constant. Any direction can be taken as the $x$ axis, so a component of $\mathbf{P}$ is conserved when the field is unchanged by arbitrary translations along that direction.

Similarly, in the cylindrical co-ordinates $r,\phi,z$, the Lagrangian is
\begin{equation*}
    L = \frac{1}{2} m (\Dot{r}^2 + r^2\Dot{\phi}^2 + \Dot{z}^2) - U(r,\phi,z),
\end{equation*}
and $\pd{L}{\Dot{\phi}} = mr^2\Dot{\phi} = M_z$ is the component of the angular momentum along the $z$ axis found in \probref{\S 9}{1}. Lagrange's equation for $\phi$ gives
\begin{equation*}
    \Dot{M}_z = \frac{d}{dt} \pd{L}{\Dot{\phi}} = \pd{L}{\phi} = -\pd{U}{\phi} .
\end{equation*}
If the field is unchanged by any rotation about the $z$ axis, then $U$ does not depend on $\phi$ and $M_z$ is constant. Here, the angular momentum is measured about the symmetry axis, \ie with the origin on that axis. The symmetries involved are continuous: the field must be unchanged by translations or rotations of any amount, however small, for the corresponding derivative of $U$ to vanish. Invariance under a few isolated rotations is not enough. The same reasoning holds for the total momentum and angular momentum of a system of particles in the field.

In each of the following figures, the body is shown on the left. The red arrows are the transformations leaving its field unchanged (a double arrow for a translation, a circular arrow for a rotation), each labelled with the quantity it conserves. The field is shown on the right in a cross-section. The arrows of the field lines give the direction of the force $\mathbf{F} = -\nabla U$ on the particle, and the dashed curves are equipotentials, along which $U$ is constant.

Let the plane be $z = 0$, as in the figure. Each point of an infinite homogeneous plane is equivalent to any other, so its field depends only on the distance to the plane, \ie $U = U(z)$. It is thus unchanged by any translation along $x$ or $y$ and by any rotation about the $z$ axis: the field lines are perpendicular to the plane and the equipotentials are planes parallel to it.
\begin{figure}[H]
    \centering
    \includegraphics{Figures/S9P3a.pdf}
\end{figure}
In terms of the force, this means that
\begin{equation*}
    F_x = -\pd{U}{x} = 0 \quad \text{and} \quad F_y = -\pd{U}{y} = 0,
\end{equation*}
and the torque about the $z$ axis vanishes:
\begin{equation*}
    \Dot{M}_z = xF_y - yF_x = 0 .
\end{equation*}
However, there is generally a force $F_z = -dU/dz$ normal to the plane, so $P_z$ is not conserved, and neither are $M_x$ and $M_y$, whose torques $yF_z$ and $-xF_z$ do not vanish. The conserved components are then
}
{
\begin{equation*}
    P_x, \quad P_y, \quad M_z
\end{equation*}
}
\end{subproblem}

\begin{subproblem}
{
That of an infinite homogeneous cylinder.
}
{
Let the axis of the cylinder be the $z$ axis, as in the figure, and $r$ the distance from that axis. The field of an infinite homogeneous cylinder is the same at every height and in every direction around the axis, so it depends only on $r$, \ie $U = U(r)$. It is thus unchanged by any translation along $z$ and by any rotation about the $z$ axis: in a cross-section, the field lines are radial and the equipotentials are circles.
\begin{figure}[H]
    \centering
    \includegraphics{Figures/S9P3b.pdf}
\end{figure}
The force $\mathbf{F} = -U'(r)\,\mathbf{\hat{r}}$ is radial. It has no axial component, so $\Dot{P}_z = F_z = 0$, and its line of action meets the axis, so it produces no torque about it. Indeed, $(F_x, F_y)$ is parallel to $(x, y)$ and
\begin{equation*}
    \Dot{M}_z = xF_y - yF_x = 0 .
\end{equation*}
The components $F_x$ and $F_y$ do not vanish, so $P_x$ and $P_y$ are not conserved, nor are $M_x$ and $M_y$. The conserved components are then
}
{
\begin{equation*}
    P_z, \quad M_z
\end{equation*}
}
\end{subproblem}

\begin{subproblem}
{
That of an infinite homogeneous prism.
}
{
Let the edges of the prism be parallel to the $z$ axis, as in the figure. The prism is the same in every cross-section perpendicular to its edges, so its field is unchanged by any translation along $z$, \ie $U = U(x,y)$, and
\begin{equation*}
    \Dot{P}_z = -\pd{U}{z} = 0 .
\end{equation*}
However, its cross-section is a polygon, so the field generally depends on the direction about any axis parallel to the edges. In the figure, the equipotentials near the prism follow the triangular cross-section and only become circular far from it. Even for a regular polygon of $n$ sides, such as the equilateral triangle shown, the field is unchanged only by the rotations by multiples of $2\pi/n$ about its axis. These isolated rotations do not make $\partial U/\partial \phi$ vanish, so they do not imply the conservation of $M_z$.
\begin{figure}[H]
    \centering
    \includegraphics{Figures/S9P3c.pdf}
\end{figure}
The only conserved component is then
}
{
\begin{equation*}
    P_z
\end{equation*}
}
\end{subproblem}

\begin{subproblem}
{
That of two points.
}
{
Let the two points be on the $z$ axis, as in the figure. The field is the sum of the fields of the two points, each depending only on the distance from the particle to that point, \ie
\begin{equation*}
    U = U_1(\rho_1) + U_2(\rho_2),
\end{equation*}
where $\rho_1$ and $\rho_2$ are the distances from the particle to the two points. A rotation of the particle about the line joining the points leaves both distances unchanged, so $U$ does not depend on $\phi$. In other words, the force exerted by each point is directed towards that point, so its line of action meets the $z$ axis and it produces no torque about it. This holds whether or not the two points have the same strength. In the figure, the field lines and equipotentials are shown in a plane containing the axis; they are the same in every such plane.
\begin{figure}[H]
    \centering
    \includegraphics{Figures/S9P3d.pdf}
\end{figure}
On the other hand, a translation moves the particle with respect to the fixed points and changes these distances, so no component of $\mathbf{P}$ is generally conserved. The only conserved component is then
}
{
\begin{equation*}
    M_z
\end{equation*}
}
\end{subproblem}

\begin{subproblem}
{
That of an infinite homogeneous half-plane.
}
{
Let the half-plane be $z = 0$, $x < 0$, as in the figure, so that its edge is the $y$ axis. A translation parallel to the edge maps the half-plane onto itself, so its field is the same in every cross-section $y = \text{constant}$, \ie $U = U(x,z)$, and
\begin{equation*}
    \Dot{P}_y = -\pd{U}{y} = 0 .
\end{equation*}
The edge removes the other symmetries of the whole plane. A translation along $x$ moves the edge, and the field depends on the distance to it: in the figure, the field lines bend towards the edge and the equipotentials curve around it. A rotation about the $z$ axis turns the edge, and a rotation about any other axis tilts the half-plane, so no rotation leaves the field unchanged and no component of $\mathbf{M}$ is generally conserved.
\begin{figure}[H]
    \centering
    \includegraphics{Figures/S9P3e.pdf}
\end{figure}
The only conserved component is then
}
{
\begin{equation*}
    P_y
\end{equation*}
}
\end{subproblem}

\begin{subproblem}
{
That of a homogeneous cone.
}
{
Let the axis of the circular cone be the $z$ axis, as in the figure. A rotation about its axis maps the cone onto itself, so its field can depend on $r$ and $z$ but not on $\phi$, \ie $U = U(r,z)$, and $M_z$ is conserved. In the figure, the field lines and equipotentials are shown in a plane containing the axis; they are the same in every such plane.
\begin{figure}[H]
    \centering
    \includegraphics{Figures/S9P3f.pdf}
\end{figure}
On the other hand, the apex of the cone is a fixed point. Any translation moves it, even one along the axis, so the translated cone is not the same and the field changes. This would remain true for an infinite cone. Hence, no component of $\mathbf{P}$ is generally conserved. The only conserved component is then
}
{
\begin{equation*}
    M_z
\end{equation*}
}
\end{subproblem}

\begin{subproblem}
{
That of a homogeneous circular torus.
}
{
Let the axis of the torus be the $z$ axis, as in the figure. A circular torus is generated by turning a circle about this axis, so any rotation about it maps the torus onto itself. Its field can then depend on $r$ and $z$ but not on $\phi$, \ie $U = U(r,z)$, and $M_z$ is conserved. In the figure, the field lines and equipotentials are shown in a plane containing the axis; they are the same in every such plane.
\begin{figure}[H]
    \centering
    \includegraphics{Figures/S9P3g.pdf}
\end{figure}
On the other hand, the torus has a fixed position, with its centre at the origin and its central circle in the plane $z = 0$. Any translation moves it to another position and changes the field, so no component of $\mathbf{P}$ is generally conserved. The only conserved component is then
}
{
\begin{equation*}
    M_z
\end{equation*}
}
\end{subproblem}

\begin{subproblem}
{
That of an infinite homogeneous cylindrical helix.
}
{
Let the axis of the helix be the $z$ axis and $R$ its radius. As in the figure, its points are
\begin{equation*}
    (x, y, z) = \left( R\cos{\phi}, R\sin{\phi}, \frac{h\phi}{2\pi} \right),
\end{equation*}
where $h$ is the pitch: the helix advances by $h$ along the axis during one positive rotation of $2\pi$. A rotation by $\delta\phi$ alone does not map the helix onto itself, nor does an arbitrary translation $\delta z$ alone (the translations by whole multiples of $h$ do, but they are isolated transformations, as the rotations of the prism). However, the combined screw displacement
\begin{equation*}
    \delta z = \frac{h}{2\pi}\delta\phi
\end{equation*}
maps each point of the helix to the point at $\phi + \delta\phi$, so it does preserve the helix and its field. Consequently, the field in the cross-section at the height $z$ is the one at $z = 0$ turned by the angle $2\pi z/h$. This is what the three cross-sections of the figure show: increasing $z$ by $h/4$ turns the pattern by $\pi/2$. In these cross-sections, the field lines show the components of the field in each plane, and not the complete three-dimensional field lines, since the field also has a component along $z$.
\begin{figure}[H]
    \centering
    \includegraphics{Figures/S9P3h.pdf}
\end{figure}
In cylindrical co-ordinates, the field thus depends on $\phi$ and $z$ only through the combination
\begin{equation*}
    \xi = z - \frac{h}{2\pi}\phi, \qquad U = U(r, \xi).
\end{equation*}
Since $\phi$ and $\phi + 2\pi$ describe the same point while $\xi$ changes by $h$, the function $U(r,\xi)$ is periodic in $\xi$ with period $h$. The Lagrangian is
\begin{equation*}
    L = \frac{1}{2} m (\Dot{r}^2 + r^2\Dot{\phi}^2 + \Dot{z}^2) - U\left(r, z - \frac{h\phi}{2\pi}\right),
\end{equation*}
and its derivatives with respect to $\phi$ and $z$ are
\begin{equation*}
    \pd{L}{\phi} = \frac{h}{2\pi} \pd{U}{\xi} \quad \text{and} \quad \pd{L}{z} = -\pd{U}{\xi}.
\end{equation*}
Using Lagrange's equations for $\phi$ and $z$, with $M_z = mr^2\Dot{\phi}$ and $P_z = m\Dot{z}$, we get
\begin{equation*}
    \Dot{M}_z = \frac{h}{2\pi} \pd{U}{\xi} \quad \text{and} \quad \Dot{P}_z = -\pd{U}{\xi},
\end{equation*}
so that
\begin{equation*}
    \frac{d}{dt} \left( M_z + \frac{h}{2\pi} P_z \right) = \frac{h}{2\pi} \pd{U}{\xi} - \frac{h}{2\pi} \pd{U}{\xi} = 0 .
\end{equation*}
No individual Cartesian component of $\mathbf{P}$ or $\mathbf{M}$ is generally conserved here. The screw symmetry conserves a combination of the axial angular momentum and the axial momentum, whose two terms have the same dimensions since $h$ is a length. The conserved quantity is then
}
{
\begin{equation*}
    M_z + \frac{h}{2\pi} P_z = \text{constant}
\end{equation*}
}
\end{subproblem}

\textbooknumber{\S 10}
\section{Mechanical similarity}

% ============ Problem 1 ============ %

\begin{problem}[title={Travel times for different masses, same path}]
{
Find the ratio of the times in the same path for particles having different masses but the same potential energy.
}
{
If the two particles have the same path, then the ratio of linear dimension is
\begin{equation} \label{S10P1_l}
    \frac{l'}{l} = \alpha = 1.
\end{equation}
We can define the ratio of time and mass by
\begin{equation*}
    \frac{t'}{t} = \beta
\end{equation*}
and
\begin{equation*}
    \frac{m'}{m} = \gamma .
\end{equation*}
Corresponding lengths scale as $\alpha$ and corresponding times as $\beta$, so the speeds at corresponding points of the two paths scale as $v'/v = \alpha/\beta$. Then, the ratio of kinetic energy is the ratio of the masses times the ratio of the squared speeds,
\begin{equation*}
    \frac{T'}{T} = \frac{m' v'^2}{m v^2} = \frac{m'}{m} \left( \frac{\alpha}{\beta} \right)^2 = \frac{\gamma \alpha^2}{\beta^2}.
\end{equation*}
To leave the equation of motion unaltered, the ratio of the kinetic energy and the potential energy must be the same, \ie,
\begin{align*}
    \frac{U'}{U} &= \frac{T'}{T} \\
    \Rightarrow \alpha^k &=  \frac{\gamma \alpha^2}{\beta^2}.
\end{align*}
Using \eqref{S10P1_l}, we get
\begin{align*}
    1 &= \frac{\gamma}{\beta^2} \\
    \Rightarrow \beta &= \sqrt{\gamma}.
\end{align*}
The ratio of the times is then
}
{
\begin{equation*}
    \frac{t'}{t} = \sqrt{\frac{m'}{m}}
\end{equation*}
}
[
We check the mass scaling by independently integrating particles of masses $m$ and $4m$ released from rest at $x=A$ in the anharmonic potential $U(x)=kx^2/2+\alpha x^4/4$, with force $F_x=-kx-\alpha x^3$. Two positive-coefficient cases are used: $(m,k,\alpha,A)=(1,1,1,1)$ and $(1.5,0.7,1.3,0.8)$ in SI units. Over $0\leq t\leq16$~s, the maximum difference $|x_{4m}(2t)-x_m(t)|$ is below $1.1\times10^{-8}$~m and the corresponding velocity difference $|2v_{4m}(2t)-v_m(t)|$ is below $1.6\times10^{-8}$~m/s. Halving the common physical RK4 step from $0.04$ to $0.02$ to $0.01$~s reduces both differences by factors greater than 16; each trajectory also converges independently to the analytic anharmonic-oscillator solution. The scaled energy drift is below $4.1\times10^{-10}$. The plots show the first case at physical and corresponding times. This is a check of selected cases, not a general proof.
\begin{figure}[H]
    \centering
    \includegraphics{Figures/S10P1-physical.pdf}
    \caption{Positions versus physical time for masses $m=1$~kg (solid) and $4m$ (dashed), released from the same position at rest in the same anharmonic well.}
    \label{fig:S10P1-physical}
\end{figure}
\begin{figure}[H]
    \centering
    \includegraphics{Figures/S10P1-rescaled.pdf}
    \caption{The independently integrated positions coincide at corresponding times: $x_m(t)$ (curve) and $x_{4m}(2t)$ (points).}
    \label{fig:S10P1-rescaled}
\end{figure}
]
\end{problem}

% ============ Problem 2 ============ %

\begin{problem}[title={Travel times for scaled potential energies}]
{
Find the ratio of the times in the same path for particles having the same mass but potential energies differing by a constant factor.
}
{
We proceed as in \probref{\S 10}{1}. If the two particles have the same path, then the ratio of linear dimension is
\begin{equation} \label{S10P2_l}
    \frac{l'}{l} = \alpha = 1.
\end{equation}
We can define the ratio of time by
\begin{equation*}
    \frac{t'}{t} = \beta .
\end{equation*}
Then, the ratio of kinetic energy is 
\begin{equation*}
    \frac{T'}{T} = \frac{\alpha^2}{\beta^2}.
\end{equation*}
The potential energy of the two particles differ by a constant factor, which mean that
\begin{equation*}
    \frac{U'}{U} = \gamma .
\end{equation*}
To leave the equation of motion unaltered, the ratio of the kinetic energy and the potential energy must be the same, \ie,
\begin{align*}
    \frac{U'}{U} &= \frac{T'}{T} \\
    \Rightarrow \gamma &=  \frac{\alpha^2}{\beta^2}.
\end{align*}
Using \eqref{S10P2_l}, we get
\begin{align*}
    \gamma &= \frac{1}{\beta^2} \\
    \Rightarrow \beta &= \sqrt{\frac{1}{\gamma}}.
\end{align*}
The ratio of the times is then
}
{
\begin{equation*}
    \frac{t'}{t} = \sqrt{\frac{U}{U'}}
\end{equation*}
}
[
We use the same two anharmonic wells and rest releases as in \probref{\S 10}{1}, keeping each mass fixed and comparing $U$ with $4U$. Both the potential and its force are multiplied by four in the second engine run. Over $0\leq t\leq16$~s, the maximum difference $|x_{4U}(t/2)-x_U(t)|$ is below $1.8\times10^{-7}$~m and $|v_{4U}(t/2)/2-v_U(t)|$ is below $2.6\times10^{-7}$~m/s. Halving the common physical RK4 step from $0.04$ to $0.02$ to $0.01$~s reduces both errors by factors greater than 16; independent comparison with the analytic anharmonic-oscillator solution and trajectory self-convergence also pass. The scaled energy drift is below $2.2\times10^{-8}$. The first case is plotted against physical and corresponding times. These are checks of selected cases, not a general proof.
\begin{figure}[H]
    \centering
    \includegraphics{Figures/S10P2-physical.pdf}
    \caption{Positions versus physical time at fixed mass $m=1$~kg for potentials $U$ (solid) and $4U$ (dashed), with the same initial position and zero initial velocity.}
    \label{fig:S10P2-physical}
\end{figure}
\begin{figure}[H]
    \centering
    \includegraphics{Figures/S10P2-rescaled.pdf}
    \caption{Positions at corresponding times: $x_U(t)$ (curve) and the independent engine samples $x_{4U}(t/2)$ (points).}
    \label{fig:S10P2-rescaled}
\end{figure}
]
\end{problem}
